% GENERATED FILE. Edit canonical lessons, then run npm run generate:books.
% canonical-source-hash: fnv1a64:ad7c2c8b674be244
% canonical-lessons: MR-C371-uttirna-hone, MR-C371-napas-hone, MR-C371-ujalni-karne, MR-C371-ghokne, MR-C371-pramanpatra

\chapter{Passing and Failing}
\label{ch:mr-passing-and-failing}
\hlchaptermodality{371}

\begin{chapteropening}
\textbf{By the end of this chapter:} \emph{I can talk about revising, passing and failing exams.}

Complete the last lesson of chapter 371: I can talk about revising, passing and failing exams.
\end{chapteropening}

\section[\texorpdfstring{\mr{उत्तीर्ण} \mr{होणे} (uttīrṇa hoṇe)}{उत्तीर्ण होणे (uttīrṇa hoṇe)}]{\mr{उत्तीर्ण} \mr{होणे} (uttīrṇa hoṇe) — to pass (an exam)}

\label{lesson:MR-C371-uttirna-hone}



\begin{quote}
Before the new one: say the Marathi for a symposium, then the Marathi for a speaker.
\end{quote}

\subsection*{You'll want to know: \mr{उत्तीर्ण} \mr{होणे}}
\textbf{\mr{उत्तीर्ण} \mr{होणे}} — \emph{uttīrṇa hoṇe} — \textquotedblleft{}to pass (an exam)\textquotedblright{}.

\begin{quote}
\mr{अरुण} \mr{पहिल्या} \mr{प्रयत्नात} \mr{परीक्षा} \mr{उत्तीर्ण} \mr{झाला}.
\end{quote}

\begin{quote}
\emph{aruṇ pahilyā prayatnāt parīkṣā uttīrṇa jhālā.}
\end{quote}

\begin{quote}
Arun passed the exam at the first attempt.
\end{quote}

Past (he …): \mr{तो} \mr{उत्तीर्ण} \mr{झाला}, \emph{to uttīrṇa jhālā}.

Formal; what is passed is the object, with no ending: \mr{परीक्षा} \mr{उत्तीर्ण} \mr{होणे}.

\begin{cousinweb}
\textbf{\mr{उत्तीर्ण} \mr{होणे}} (\emph{uttīrṇa hoṇe}) joins the tatsama \mr{उत्तीर्ण}, Sanskrit \emph{uttīrṇa}, \textquotedblleft{}crossed over, landed\textquotedblright{} (Monier-Williams), to \mr{होणे}, \textquotedblleft{}to become\textquotedblright{}: to get across.
\end{cousinweb}

\begin{grammarlens}[title={What you've built}]
One more word for giving reasons and opinions.
\end{grammarlens}

\subsection*{Your turn}
\begin{itemize}
\raggedright
  \item \emph{Say it:} \emph{uttīrṇa hoṇe}
  \item \emph{Say it:} \emph{uttīrṇa hoṇe}, once more
  \item \emph{Say it:} \emph{aruṇ pahilyā prayatnāt parīkṣā uttīrṇa jhālā.} — Arun passed the exam at the first attempt.
  \item \emph{Recall:} say the Marathi for a symposium, then the Marathi for a speaker, then say \emph{uttīrṇa hoṇe} again
\end{itemize}

\begin{tcolorbox}[breakable,colback=teal!4,colframe=teal!35!black,title={Before you move on}]
What does \mr{उत्तीर्ण} \mr{होणे} mean? (\textquotedblleft{}To pass (an exam)\textquotedblright{}.) Say it once more.
\end{tcolorbox}
\section[\texorpdfstring{\mr{नापास} \mr{होणे} (nāpās hoṇe)}{नापास होणे (nāpās hoṇe)}]{\mr{नापास} \mr{होणे} (nāpās hoṇe) — to fail (an exam)}

\label{lesson:MR-C371-napas-hone}



\begin{quote}
Before the new one: say the Marathi for a speaker, then the Marathi for to pass.
\end{quote}

\subsection*{You'll want to know: \mr{नापास} \mr{होणे}}
\textbf{\mr{नापास} \mr{होणे}} — \emph{nāpās hoṇe} — \textquotedblleft{}to fail (an exam)\textquotedblright{}.

\begin{quote}
\mr{तो} \mr{गणितात} \mr{नापास} \mr{झाला}.
\end{quote}

\begin{quote}
\emph{to gaṇitāt nāpās jhālā.}
\end{quote}

\begin{quote}
He failed in maths.
\end{quote}

Past (he …): \mr{तो} \mr{नापास} \mr{झाला}, \emph{to nāpās jhālā}.

The school subject takes -\mr{त}: \mr{गणितात} \mr{नापास}, failed in maths.

\begin{cousinweb}
\textbf{\mr{नापास} \mr{होणे}} (\emph{nāpās hoṇe}) joins \mr{नापास}, a hybrid of the Persian prefix \emph{nā-}, \textquotedblleft{}not\textquotedblright{}, and a loan from English \emph{pass}, to \mr{होणे}, \textquotedblleft{}to become\textquotedblright{}.
\end{cousinweb}

\begin{grammarlens}[title={What you've built}]
One more word for giving reasons and opinions.
\end{grammarlens}

\subsection*{Your turn}
\begin{itemize}
\raggedright
  \item \emph{Say it:} \emph{nāpās hoṇe}
  \item \emph{Say it:} \emph{nāpās hoṇe}, once more
  \item \emph{Say it:} \emph{to gaṇitāt nāpās jhālā.} — He failed in maths.
  \item \emph{Recall:} say the Marathi for a speaker, then the Marathi for to pass, then say \emph{nāpās hoṇe} again
\end{itemize}

\begin{tcolorbox}[breakable,colback=teal!4,colframe=teal!35!black,title={Before you move on}]
What does \mr{नापास} \mr{होणे} mean? (\textquotedblleft{}To fail (an exam)\textquotedblright{}.) Say it once more.
\end{tcolorbox}
\section[\texorpdfstring{\mr{उजळणी} \mr{करणे} (ujaḷṇī karṇe)}{उजळणी करणे (ujaḷṇī karṇe)}]{\mr{उजळणी} \mr{करणे} (ujaḷṇī karṇe) — to revise}

\label{lesson:MR-C371-ujalni-karne}



\begin{quote}
Before the new one: say the Marathi for to pass, then the Marathi for to fail.
\end{quote}

\subsection*{You'll want to know: \mr{उजळणी} \mr{करणे}}
\textbf{\mr{उजळणी} \mr{करणे}} — \emph{ujaḷṇī karṇe} — \textquotedblleft{}to revise\textquotedblright{}.

\begin{quote}
\mr{परीक्षेआधी} \mr{मी} \mr{सगळ्या} \mr{धड्यांची} \mr{उजळणी} \mr{केली}.
\end{quote}

\begin{quote}
\emph{parīkṣeādhī mī sagaḷyā dhaḍyāñcī ujaḷṇī kelī.}
\end{quote}

\begin{quote}
Before the exam I revised all the lessons.
\end{quote}

Past, with -\mr{ने} (he …): \mr{त्याने} \mr{उजळणी} \mr{केली}, \emph{tyāne ujaḷṇī kelī}.

What is revised takes -\mr{ची}: \mr{धड्यांची} \mr{उजळणी}.

\begin{cousinweb}
\textbf{\mr{उजळणी} \mr{करणे}} (\emph{ujaḷṇī karṇe}) joins \mr{उजळणी}, \textquotedblleft{}brightening up again\textquotedblright{}, built on the Marathi verb \mr{उजळणे}, \textquotedblleft{}to polish, kindle\textquotedblright{}, inherited from Sanskrit \emph{ujjvalati}, \textquotedblleft{}blazes up\textquotedblright{} (CDIAL 1671), to \mr{करणे}.
\end{cousinweb}

\begin{grammarlens}[title={What you've built}]
One more word for giving reasons and opinions.
\end{grammarlens}

\subsection*{Your turn}
\begin{itemize}
\raggedright
  \item \emph{Say it:} \emph{ujaḷṇī karṇe}
  \item \emph{Say it:} \emph{ujaḷṇī karṇe}, once more
  \item \emph{Say it:} \emph{parīkṣeādhī mī sagaḷyā dhaḍyāñcī ujaḷṇī kelī.} — Before the exam I revised all the lessons.
  \item \emph{Recall:} say the Marathi for to pass, then the Marathi for to fail, then say \emph{ujaḷṇī karṇe} again
\end{itemize}

\begin{tcolorbox}[breakable,colback=teal!4,colframe=teal!35!black,title={Before you move on}]
What does \mr{उजळणी} \mr{करणे} mean? (\textquotedblleft{}To revise\textquotedblright{}.) Say it once more.
\end{tcolorbox}
\section[\texorpdfstring{\mr{घोकणे} (ghokṇe)}{घोकणे (ghokṇe)}]{\mr{घोकणे} (ghokṇe) — to cram; to learn by rote}

\label{lesson:MR-C371-ghokne}



\begin{quote}
Before the new one: say the Marathi for to fail, then the Marathi for to revise.
\end{quote}

\subsection*{You'll want to know: \mr{घोकणे}}
\textbf{\mr{घोकणे}} — \emph{ghokṇe} — \textquotedblleft{}to cram; to learn by rote\textquotedblright{}.

\begin{quote}
\mr{न} \mr{समजता} \mr{उत्तरं} \mr{घोकून} \mr{काही} \mr{उपयोग} \mr{नाही}.
\end{quote}

\begin{quote}
\emph{na samajtā uttaraṁ ghokūn kāhī upyog nāhī.}
\end{quote}

\begin{quote}
There is no use cramming answers without understanding them.
\end{quote}

Past, with -\mr{ने} (he …): \mr{त्याने} \mr{घोकलं}, \emph{tyāne ghoklaṁ}.

\mr{घोकंपट्टी}, rote learning.

\begin{cousinweb}
\textbf{\mr{घोकणे}} (\emph{ghokṇe}) is a Marathi verb whose origin is not settled.
\end{cousinweb}

\begin{grammarlens}[title={What you've built}]
One more word for giving reasons and opinions.
\end{grammarlens}

\subsection*{Your turn}
\begin{itemize}
\raggedright
  \item \emph{Say it:} \emph{ghokṇe}
  \item \emph{Say it:} \emph{ghokṇe}, once more
  \item \emph{Say it:} \emph{na samajtā uttaraṁ ghokūn kāhī upyog nāhī.} — There is no use cramming answers without understanding them.
  \item \emph{Recall:} say the Marathi for to fail, then the Marathi for to revise, then say \emph{ghokṇe} again
\end{itemize}

\begin{tcolorbox}[breakable,colback=teal!4,colframe=teal!35!black,title={Before you move on}]
What does \mr{घोकणे} mean? (\textquotedblleft{}To cram; to learn by rote\textquotedblright{}.) Say it once more.
\end{tcolorbox}
\section[\texorpdfstring{\mr{प्रमाणपत्र} (pramāṇpatra)}{प्रमाणपत्र (pramāṇpatra)}]{\mr{प्रमाणपत्र} (pramāṇpatra) — a certificate}

\label{lesson:MR-C371-pramanpatra}



\begin{quote}
Before the new one: say the Marathi for to revise, then the Marathi for to cram.
\end{quote}

\subsection*{You'll want to know: \mr{प्रमाणपत्र}}
\textbf{\mr{प्रमाणपत्र}} — \emph{pramāṇpatra} — \textquotedblleft{}a certificate\textquotedblright{}.

\begin{quote}
\mr{स्पर्धेनंतर} \mr{सगळ्यांना} \mr{प्रमाणपत्र} \mr{मिळालं}.
\end{quote}

\begin{quote}
\emph{spardhenantar sagaḷyānnā pramāṇpatra miḷālaṁ.}
\end{quote}

\begin{quote}
After the competition everyone got a certificate.
\end{quote}

Neuter.

Literally a letter of proof.

\begin{cousinweb}
\textbf{\mr{प्रमाणपत्र}} (\emph{pramāṇpatra}) is a learned compound built from Sanskrit \emph{pramāṇa}, \textquotedblleft{}measure, proof\textquotedblright{} (CDIAL 8723), and \emph{patra}, \textquotedblleft{}leaf, letter\textquotedblright{}.
\end{cousinweb}

\begin{grammarlens}[title={What you've built}]
One more word for giving reasons and opinions.
\end{grammarlens}

\subsection*{Your turn}
\begin{itemize}
\raggedright
  \item \emph{Say it:} \emph{pramāṇpatra}
  \item \emph{Say it:} \emph{pramāṇpatra}, once more
  \item \emph{Say it:} \emph{spardhenantar sagaḷyānnā pramāṇpatra miḷālaṁ.} — After the competition everyone got a certificate.
  \item \emph{Recall:} say the Marathi for to revise, then the Marathi for to cram, then say \emph{pramāṇpatra} again
\end{itemize}

\begin{tcolorbox}[breakable,colback=teal!4,colframe=teal!35!black,title={Before you move on}]
What does \mr{प्रमाणपत्र} mean? (\textquotedblleft{}A certificate\textquotedblright{}.) Say it once more.
\end{tcolorbox}
